\exercisesection{§ 3 的习题}

\begin{enumerate}
\def\labelenumi{\arabic{enumi})}
\item
  记号与假设均如 no. 1 所述，证明相对于 W 的诸室是开单纯形当且仅当 W 无限且不可约。证明 E/W 是紧的当且仅当 W 是一些无限不可约群的乘积。
\item
  设 V 是一个赋有内积的有限维实向量空间，F 是 V 的正交群中由反射生成的有限子群，\(\Lambda\) 是 V 的一个在 F 下稳定的离散子群，W 是 V 中由 F 与以属于 \(\Lambda\) 的向量为平移向量的诸平移所生成的位移群。设 H 是 V 中使 \(s_{H} \in F\) 的超平面 H 的集合。设 R 是与 H 的某个元素正交的 \(\Lambda\) 中元素的集合。
\end{enumerate}

\begin{enumerate}
\def\labelenumi{\alph{enumi})}
\item
  W 由反射生成当且仅当 R 生成 \(\Lambda\)。
\item
  在 \(\mathbf{R}^2\)（赋有内积 \(((x,y),(x',y'))\mapsto xx' + yy'\)）中，设
\end{enumerate}

\[
e _ {1} = (1, 0), \quad e _ {2} = \left(- \frac {1}{2}, \frac {\sqrt {3}}{2}\right), \quad e _ {3} = \left(- \frac {1}{2}, - \frac {\sqrt {3}}{2}\right), \quad \Delta_ {i} = \mathbf {R} e _ {i},
\]

F 是由诸 \(s_{\Delta_{i}}\) 生成的二面体群，\(\Lambda\) 是 \(R^{2}\) 的一个离散子群，由诸 \(e_{i}\) 生成，该子群在 F 下稳定。证明 W 不由反射生成。

\begin{enumerate}
\def\labelenumi{\arabic{enumi})}
\setcounter{enumi}{2}
\item
  设 V 是一个有限维实向量空间，W 是 \(\mathbf{GL}(V)\) 的一个由反射生成的有限子群。证明 W 的每个 2 阶元素都是 W 中两两交换的反射的乘积。（对 \(\dim(V)\) 用归纳法论证，并使用 Prop. 2。）
\item
  设 V 是一个有限维实向量空间，W 是 \(\mathbf{GL}(V)\) 的一个由反射生成的有限子群，w 是 \(W\) 的一个元素，\(V'\) 是 V 的一个在 w 下稳定的向量子空间，k 是 \(w|V'\)（即 w 到 \(V'\) 上的限制）的阶。证明存在 \(x \in W\)，其阶为 k，使 \(V'\) 稳定，且 \(x|V' = w|V'\)。（设 \(W'\) 是 W 中使 \(V'\) 的点保持不动的元素所成的集合。这个群由反射生成，且 w 置换相对于 \(W'\) 的诸室；由此推出存在 \(h \in W'\) 使得 wh 使 \(W'\) 的一个室稳定；证明可取 x = wh。）
\end{enumerate}

¶ 5) 设 V 是一个有限维实向量空间，W 是 GL(V) 的一个由反射生成的有限子群，C 是 W 的一个室，S 是 C 的墙的集合。若 J ⊂ S，设 W\_J 是 W 中由 H ∈ J 的诸 s\_H 生成的子群，并设 \(\varepsilon(J) = (-1)^{\operatorname{Card}(J)}\)。证明

\[
\frac {1}{\operatorname{Card} (\mathrm{W})} = \sum_ {\mathrm{J} \subset \mathrm{S}} \frac {\varepsilon (\mathrm{J})}{\operatorname{Card} (\mathrm{W} _ {\mathrm{J}})}.\tag{(*)}
\]

（设 \((x|y)\) 是 V 上在 W 下不变的内积，\(\Sigma\) 是 V 的单位球面，\(\mu\) 是 \(\Sigma\) 上在 W 下不变且总质量为 1 的正测度。若 \(H \in S\)，设 \(D_{H}\) 是由 H 决定且包含 C 的开半空间。设 \(E_{H} = D_{H} \cap \Sigma\)。则

\[
(- \overline {{C}}) \cap \Sigma = \bigcap_ {H \in S} (\Sigma - E _ {H}),
\]

于是

\[
\begin{array}{l} \frac {1}{\operatorname{Card} (\mathrm{W})} = \mu (\overline {{C}} \cap \Sigma) = \int_ {\Sigma} \prod_ {\mathrm{H} \in \mathrm{S}} (1 - \varphi_ {\mathrm{E} _ {\mathrm{H}}}) d \mu \\ = \sum_ {\mathrm{J} \subset \mathrm{S}} \varepsilon (\mathrm{J}) \mu \left(\bigcap_ {\mathrm{H} \in \mathrm{J}} \mathrm{E} _ {\mathrm{H}}\right). \end{array}
\]

最后注意 \(\bigcap_{H\in J}E_H\) 是 \(\Sigma\) 与 \(W_J\) 的一个室的交，故其测度等于 \(1 / \mathrm{Card}(W_J)\)，由此即得结论。）

借助 Chap. IV, §1 的 Exerc. 26 e) 重新得到公式 (*)（在所考虑的习题中证明的恒等式里取 t = 1）。

\begin{enumerate}
\def\labelenumi{\arabic{enumi})}
\setcounter{enumi}{5}
\tightlist
\item
  \begin{enumerate}
  \def\labelenumii{\alph{enumii})}
  \tightlist
  \item
    设 K 是一个交换域，V 是 K 上维数为有限 n 的向量空间，\(\varphi\) 是 V 上一个对称双线性型，N 是 \(\varphi\) 的核。假设 \(\dim N = 1\)。证明 \(\varphi\) 到 \(\bigwedge^{n-1} V\) 上的扩张的核的维数为 n - 1。
  \end{enumerate}
\end{enumerate}

\begin{enumerate}
\def\labelenumi{\alph{enumi})}
\setcounter{enumi}{1}
\tightlist
\item
  再假设 K = R 且 \(\varphi\) 是正的。设 \((e_{1},\ldots,e_{n})\) 是 V 的基，\(a_{ij} = \varphi(e_{i},e_{j})\)。假设 \(a_{ij} \leqslant 0\)（对 \(i \neq j\)）。假设不存在 \(\{1,2,\ldots,n\}\) 的分割 I ∪ J 使得 \(a_{ij} = 0\)（对 \(i \in I, j \in J\)）。设 \(A_{ij}\) 是 \(a_{ij}\) 在矩阵 \((a_{ij})\) 中的余子式。证明对所有 i 与 j 有 \(A_{ij} > 0\)。（设 \(\zeta_{1}e_{1} + \cdots + \zeta_{n}e_{n}\) 是一个生成 N 的向量，其所有坐标都 \textgreater{} 0。证明 \((A_{ij})\) 的每一行与每一列都与 \((\zeta_{1},\ldots,\zeta_{n})\) 成比例。
\end{enumerate}

由此推出 \(\mathbf{A}_{ij} = \mu \zeta_i\zeta_j\)（对某个常数 \(\mu\)）。通过考察诸 \(\mathbf{A}_{ii}\)，证明 \(\mu > 0\)。）

\begin{enumerate}
\def\labelenumi{\alph{enumi})}
\setcounter{enumi}{2}
\tightlist
\item
  证明 \(\zeta_1, \ldots, \zeta_n\) 与 \(\sqrt{\mathrm{A}_{11}}, \ldots, \sqrt{\mathrm{A}_{nn}}\) 成比例。
\end{enumerate}

\begin{enumerate}
\def\labelenumi{\arabic{enumi})}
\setcounter{enumi}{6}
\tightlist
\item
  设 \(q(\xi_1, \ldots, \xi_n) = \sum_{i,j} a_{ij} \xi_i \xi_j\)（\(a_{ij} = a_{ji}\)）是 \(\mathbf{R}^n\) 上一个退化的正二次型，满足 \(a_{ij} \leqslant 0\)（对 \(i \neq j\)）。假设不存在分割 \(I \cup J\)（\(\{1, 2, \ldots, n\}\) 的分割）使得 \(a_{ij} = 0\)（对 \(i \in I, j \in J\)）。
\end{enumerate}

\begin{enumerate}
\def\labelenumi{\alph{enumi})}
\item
  证明：若置 \(\xi_{i}=0\)，则得到关于 \(\xi_{1},\ldots,\xi_{i-1},\xi_{i+1},\ldots,\xi_{n}\) 的一个非退化正二次型。
\item
  证明对所有 i 有 \(a_{ii} > 0\)。（设 \((\zeta_{1}, \ldots, \zeta_{n})\) 是 q 的核中的一个元素，满足 \(\zeta_{1} > 0, \ldots, \zeta_{n} > 0\)。利用等式 \(a_{1i}\zeta_{1} + \cdots + a_{ni}\zeta_{n} = 0\)。）
\item
  证明：若把某个 \(a_{ij}\) 换成某个 \(a_{ij}' < a_{ij}\)，则新的二次型不是正的。（利用等式 \(\sum_{i,j} a_{ij} \zeta_i \zeta_j = 0\)。）
\end{enumerate}

\begin{enumerate}
\def\labelenumi{\arabic{enumi})}
\setcounter{enumi}{7}
\tightlist
\item
  设 \((a_{ij})\) 是一个有 \(n\) 行 \(n\) 列的实对称矩阵。
\end{enumerate}

\begin{enumerate}
\def\labelenumi{\alph{enumi})}
\tightlist
\item
  置 \(s_k = \sum_{i=1}^{n} a_{ik}\)。则对所有 \(\xi_1, \ldots, \xi_n \in \mathbf{R}\)，
\end{enumerate}

\[
\sum_ {i, k} a _ {i k} \xi_ {i} \xi_ {k} = \sum_ {k} s _ {k} \xi_ {k} ^ {2} - \frac {1}{2} \sum_ {i, k} a _ {i k} (\xi_ {i} - \xi_ {k}) ^ {2}.
\]

\begin{enumerate}
\def\labelenumi{\alph{enumi})}
\setcounter{enumi}{1}
\tightlist
\item
  设 \(\zeta_1, \ldots, \zeta_n \in \mathbf{R}^*\)。置 \(\sum_{i} \zeta_i a_{ik} = t_k\)。则
\end{enumerate}

\[
\sum_ {i, k} a _ {i k} \xi_ {i} \xi_ {k} = \sum_ {k} \frac {t _ {k} \xi_ {k} ^ {2}}{\zeta_ {k}} - \frac {1}{2} \sum_ {i, k} \zeta_ {i} \zeta_ {k} a _ {i k} \left(\frac {\xi_ {i}}{\zeta_ {i}} - \frac {\xi_ {k}}{\zeta_ {k}}\right) ^ {2}.
\]

（在 \(a\) ) 中，把 \(\xi_{i}\) 换成 \(\frac{\xi_i}{\zeta_i}\)，把 \(a_{ik}\) 换成 \(\zeta_i\zeta_k a_{ik}\)。）

\begin{enumerate}
\def\labelenumi{\alph{enumi})}
\setcounter{enumi}{2}
\item
  若存在数 \(\zeta_1, \ldots, \zeta_n > 0\) 使得 \(\sum_{i} \zeta_i a_{ik} = 0\)（\(k = 1, 2, \ldots, n\)），且 \(a_{ij} \leqslant 0\)（对 \(i \neq j\)），则二次型 \(\sum_{i,k} a_{ik} \xi_i \xi_k\) 是退化且正的。（利用 b)。）
\item
  设 \(\sum_{i,j} q_{ij} \xi_i \xi_j\) 是 \(\mathbf{R}^n\) 上一个二次型，满足 \(q_{ij} \leqslant 0\)（对 \(i \neq j\)）。假设不存在分割 \(\mathrm{I} \cup \mathrm{J}\)（\(\{1,2,\ldots,n\}\) 的分割）使得 \(q_{ij} = 0\)（对 \(i \in \mathrm{I}, j \in \mathrm{J}\)）。则该二次型非退化且正当且仅当存在 \(\zeta_1 > 0, \ldots, \zeta_n > 0\) 使得 \(\sum_{j} \zeta_i q_{ik} = 0 (k = 1, \ldots, n)\)。
\end{enumerate}

